


\def\stat{ilini}





\def\tit{МЕТОД ЦЕЛЕВОГО ПЕРЕМЕЩЕНИЯ РЕШЕНИЯ В~ТЕХНОЛОГИЯХ СИТУАЦИОННОГО 


УПРАВЛЕНИЯ}





\def\titkol{Метод целевого перемещения решения в~технологиях ситуационного 


управления}





\def\aut{А.\,В.~Ильин$^1$, В.\,Д.~Ильин$^2$}





\def\autkol{А.\,В.~Ильин, В.\,Д.~Ильин}





\titel{\tit}{\aut}{\autkol}{\titkol}





\index{Ильин А.\,В.}


\index{Ильин В.\,Д.}


\index{Ilyin A.\,V.}


\index{Ilyin V.\,D.}





%{\renewcommand{\thefootnote}{\fnsymbol{footnote}}


%\footnotetext[1]


%{Результаты 


%получены при выполнении на\-уч\-но-ис\-сле\-до\-ва\-тель\-ской работы 


%<<Моделирование социальных, экономических и~экологических 


%процессов>> (№\,0063-2016-0005) в~соответствии 


%с~государственным заданием ФАНО России для Федерального 


%исследовательского центра <<Информатика и~управ\-ле\-ние>> РАН.}} 





\renewcommand{\thefootnote}{\arabic{footnote}}


\footnotetext[1]{Государственный научно-исследовательской институт авиационных систем,\\ \mbox{ilyin@res-plan.com}}


\footnotetext[2]{Федеральный исследовательский центр <<Информатика и~управ\-ле\-ние>> Российской академии наук, 


\mbox{vdilyin@yandex.ru}}





 \label{st\stat}





 \thispagestyle{headings}


 


 \vspace*{-10pt} 


  











\Abst{Рассмотрена обновленная технология интерактивного решения линейных задач 


ситуационного управления методом целевого перемещения решения. Формулировки задач


 и~требования к~качеству решения формируются на основе анализа портретов целевой 


 и~достигнутой ситуации. Метод целевого перемещения решения (в~отличие от методов 


линейного программирования) позволяет получить требуемый результат и~при 


несовместности системы ограничений. Интерактивность метода предоставляет 


возможность эксперту на каждом шаге поиска решения изменять формулировку задачи, 


исходя из его представлений о~реализуемости и~эффективности. Технология 


ориентирована на реализацию в~виде он\-лайн-сер\-ви\-сов.} 





\KW{ситуационное управление; линейные задачи ситуационного управ\-ле\-ния; 


распределение ресурсов; портрет ситуации; метод целевого перемещения решения;  


он\-лайн-сер\-ви\-сы}





\DOI{10.14357\!/\!08696527220412}  





\vskip 10pt plus 9pt minus 6pt





 \thispagestyle{headings}


 


 \vspace*{-16pt}





\section{Введение}





\vspace*{-2pt}





  В наши дни интенсивного развития облачных вычислений  


и~ин\-тер\-нет-сер\-ви\-сов различного назначения (навигационных, 


образовательных и~др.) внимание исследователей и~разработчиков 


информационных технологий привлечено к~идее комплексной автоматизации 


различных видов деятельности на основе ин\-тер\-нет-сер\-ви\-сов. С~середины 


1990-х~гг.\ значение успешной реализации этой идеи [получившей название 


\textit{digital economy} (цифровая экономика)] неуклонно рас\-тет 


и~связывается с~конкурентоспособностью корпораций и~стран~[1--3]. 


В~2017~г.\ в~нашей стране утверждена программа <<Цифровая 


экономика Российской Федерации>>~[4]. 


  


  Увеличиваются темпы применения технологий M2M  


(Machine-to-Machine)~[5, 6], IoT (Internet of Things)~[7, 8] и~IoS (Internet of 


Services)~[9--12], которые становятся эффективным средством циф\-ро\-ви\-за\-ции 


различных областей де\-я\-тель\-ности. Растет число M2M-, IoT- и~IoS-применений 


в~энергетике, промышленности, сельском хозяйстве, здравоохранении и~других 


сферах деятельности. Темпы распространения технологий M2M, IoT и~IoS во 


многом определяются снижением сто\-и\-мости вы\-чис\-ли\-тел\-ьных устройств 


и~передачи данных, рос\-том до\-ступ\-ности гиб\-ких сис\-тем хранения и~анализа 


данных (благодаря развитию <<облачных>> технологий и~технологий 


<<больших данных>>). 


  


  Постоянно актуальной остается проблема методологического обеспечения 


технологий решения задач ситуационного управ\-ле\-ния слож\-ны\-ми  


че\-ло\-ве\-ко-ма\-шин\-ны\-ми сис\-те\-ма\-ми~\cite{13-il, 14-il, 15-il, 16-il}.


  


  \textbf{О применимости методов линейного программирования.} 


Применимость методов линейного программирования (ЛП)~\cite{17-il} для 


решения линейных задач ситуационного управления существенно ограничена 


предположением о~том, что прогнозируемые входные данные, используемые 


для поиска решения, совпадают с~реальными при реализации решения. 


Известно, что в~практически значимых ЛП-за\-да\-чах решений, как правило, не 


существует из-за несовместности системы ограничений. Кроме того, расчетные 


входные данные могут значительно отличаться от реальных в~процессе 


реализации решения. В~предельных случаях за время между получением 


решения и~его реализацией могут измениться множество основных 


переменных и~система ограничений. 


  


  Обычно ЛП-задачи решаются сим\-плекс-ме\-то\-дом~\cite{17-il} или 


методом внутренних точек~\cite{18-il}. Решения могут быть найдены только 


при совместности системы ограничений. В~случае несовместности может быть 


решена задача поиска чебышёвской точки, которая также решается симплекс-ме\-то\-дом. 


Известно, что, во-пер\-вых, ЛП-за\-да\-ча относится к~числу 


математически некорректных (из-за неустойчивости решения к~малым 


изменениям входных данных)~\cite{19-il}; во-вто\-рых, решения, найденные 


как чебышёвская точка, неприменимы, так как нарушают заданные 


ограничения.


  


  Чтобы создать метод решения практически значимых линейных задач, была 


предложена идея неформальной постановки таких задач на основе 


\textit{систем обязательных и~ориентирующих требований}~\cite{16-il, 20-il}. 


Идея была реализована в~\textit{методе целевого перемещения 


решения}~\cite{20-il, 21-il, 22-il, 23-il, 24-il}. 


  


  





\section{Портреты ситуаций: формирование и~анализ }


  


  \textit{Целевой} будем называть ситуацию, которую планируется создать 


в~результате управляющего воздействия на объект ситуационного управ\-ле\-ния; 


\textit{стартовой}~--- применительно к~которой проектируется управ\-ля\-ющее 


воздействие; \textit{достигнутой}~--- полученную в~результате управ\-ля\-юще\-го 


воздействия.


  


  \textit{Портрет ситуации}~--- описание, включающее данные о~со\-сто\-янии 


объекта, располагаемых вариантах управляющих воздействий и~ресурсного 


обеспечения их реализации~\cite{16-il}.


  


  Комплекс формирования и~анализа портретов ситуаций работает 


  в~соответствии с~\textit{системой правил} $\langle R, S\rangle$, где $R$~--- множество 


правил; $S\hm\subseteq R\times R$~--- бинарное отношение на~$R$. 


\textit{Правило} имеет вид либо $D_1\Rightarrow D_2$, либо $D\hm\Rightarrow 


I$ (здесь $D$, $D_1$ и~$D_2$~--- произвольные описания; $I$~--- любая инструкция). Правила 


упорядочены с~помощью отношения \textit{специализации} $S$, 


замыкание~$S^\circ$ которого иррефлексивно (т.\,е.\ ни для какого $r\hm\in R$ 


не имеет места $\langle r, r\rangle \hm \in S^\circ$; замыкание отношения $S$~--- 


множество $S^\circ$ пар $\langle r, r^\prime \rangle$ таких, что $\langle r^\prime, 


r^{\prime\prime}\rangle \hm\in S^\circ$ в~том и~только в~том случае, когда 


существует последовательность правил $r^\prime\hm= r_0, r_1,\ldots ,r_m \hm= 


r^{\prime\prime}$ ($m\hm \geq 1$), причем $\langle r_i, r_{i+1}\rangle \hm\in S$ для каждого 


$0 \hm\leq  i \hm\leq m\hm- 1$). Если для любых двух правил $r_1$ и~$r_2$ указано, 


что $r_2$ является \textit{специализацией}~$r_1$, то~$r_1$ является 


\textit{обобщением} $r_2$. Каждое \textit{пра\-ви\-ло-спе\-ци\-а\-ли\-за\-ция} 


предназначено для срабатывания в~более част\-ных ситуациях, чем 


со\-от\-вет\-ст\-ву\-ющее  


\textit{пра\-ви\-ло-обоб\-ще\-ние}. Любое пра\-ви\-ло-спе\-ци\-а\-ли\-за\-ция 


имеет приоритет в~срабатывании перед со\-от\-вет\-ст\-ву\-ющим  


пра\-ви\-лом-обоб\-ще\-нием.


  


  \textit{Механизм интерпретации} запроса на множестве правил указывает, 


какие правила и~в~каком порядке применять, а~так\-же какие описания считать 


истинными. За\-прос пред\-став\-ля\-ет собой набор данных, который включает 


значения параметров пространства со\-сто\-яний объекта, спецификации 


располагаемых типов управляющих воздействий и~ресурсов, необходимых для 


реализации воздействий~\cite{21-il}. 





\section{Общая линейная задача распределения ресурсов, решаемая~методом~целевого~перемещения решения}


  


  Пусть $a_{ij}$ ($i\hm=1,\ldots , m$, $j\hm=1,\ldots , n$)~--- расход $i$-го ресурса 


при единичной интенсивности использования $j$-го варианта (средства) 


деятельности; $b_i$ ($i\hm=1,\ldots , m$)~--- рас\-по\-ла\-га\-емый объем $i$-го ресурса; 


$x_j$ ($j\hm=1,\ldots , n$)~--- искомая ин\-тен\-сив\-ность использования $j$-го 


средства. Суммарный расход $i$-го ресурса выражается линейной формой 


$a_{i1}x_1+ \cdots + a_{in}x_n$. Таким образом, составное требование 


к~искомому распределению имеет вид $\{a_{i1}x_1 +\cdots + a_{in}x_n \hm\leq bi\}$ 


($i\hm=1,\ldots , m$). Кроме того, может быть задан набор показателей 


эф\-фек\-тив\-ности общей де\-я\-тель\-ности: $\{c_{i1}x_1 +\cdots +c_{in}x_n\}$ 


($i\hm=1,\ldots , k$), где $c_{ij}$~--- удельный показатель полезности $i$-го типа от 


использования единицы интенсивности действий $j$-го средства. Для всякого 


показателя эф\-фек\-тив\-ности может быть задано некоторое прос\-тое правило. 


Также могут быть заданы приоритеты требований~$p_i$ ($0\hm< p_i \hm\leq \infty$, 


$i\hm\in  [1, m\hm+k]$).


  


  В общем случае для каждой ресурсной функ\-ции может быть определено 


двустороннее ограничение (конъюнкция двух простых требований). Таким 


образом, общую сис\-те\-му требований мож\-но записать в~сле\-ду\-ющем виде:


  $$


\left\{ \left[b_i\leq \right] a_{i1}x_1 +\cdots + a_{in}x_n \left[\leq B_i\right]\left[p_i\right];\ x_j 


\geq0\right\}\ (i = 1,\ldots , m + k,\ j=1,\ldots , n)


$$


(все коэффициенты ресурсных ограничений и~показателей эф\-фек\-тив\-ности 


обозначены через $a_{ij}$; квадратные скоб\-ки обозначают не\-обя\-за\-тель\-ность 


задания ограничений и~приоритетов; переменные~$x_j$ неотрицательны 


в~соответствии со смыс\-лом задачи).


  


  \textit{Простое требование} может быть пред\-став\-ле\-но одной из трех форм: 


  $$


  F_i(\bm{x})=c_i[p_i];\enskip F_i(\bm{x}) \leq c_i[p_i];\enskip 


F_i(\bm{x}) \geq c_i[p_i],


$$


 где $F_i(\bm{x})$~--- $i$-я функция (например, 


ресурсная или оценочная); $\bm{x}$~--- вектор решения (здесь и~далее векторы 


и~мат\-ри\-цы выделены жир\-ным шриф\-том); $c_i$~--- константа; $p_i$~--- 


приоритет требования ($0\hm<p_i \hm\leq \infty$) (квадратные скобки обозначают 


необязательность задания приоритета). \textit{Составное требование}~--- 


логическая комбинация простых требований, полученная с~помощью 


логических связок конъюнкции, дизъюнкции и~отрицания ($\wedge$, $\vee$ и~$\neg$). 


\textit{Обязательные требования} не могут быть нарушены (имеют 


абсолютный приоритет $p_i\hm=\infty$). \textit{Ориентирующие требования} 


задают предпочтительные значения функций ($p_i\hm > 0$).


  


  Задача заключается в~отыскании вектора распределения ресурсов 


$\bm{x}\hm= (x_1 \cdots x_n)$, которому соответствуют значения ре\-сурс\-ных 


функций, оцениваемые экспертом как наиболее эффективные в~данной 


ситуации. 


  


  В процессе поиска решения эксперт решает некоторую по\-сле\-до\-ва\-тель\-ность 


част\-ных задач, име\-ющих формальные по\-ста\-нов\-ки и~алгоритмы, и~выполняет 


сравнительный анализ решений~\cite{24-il}.


  


  Начальное решение (начальная точка) может быть выбрано экспертом 


произвольно. По умолчанию предлагается компромиссное решение 


(чебышёвская точка). На любом шаге поиска решения эксперт, 


проанализировав значения функций в~составе требований и~оценив степень 


эффективности и~реализуемости найденного решения, может изменить любые 


требования. В~частности, может зафиксировать значение некоторой функции 


(если оно устраивает его). Таким образом эксперт пошагово приближается 


к~решению, по его мнению, со\-че\-та\-юще\-му свойства эффективности 


и~реализуемости. Любое решение может быть занесено в~базу возможных 


решений для последующего анализа.


  


  Шаг целевого перемещения решения выполняется следующим способом. 


Пусть $\bm{x}^\prime \hm = (x_1^\prime\cdots  x_n^\prime)$ ($x_j^\prime \hm\geq 0$, $ j \hm= 


1,\ldots , n$)~--- решение, полученное на предыду\-щем шаге, и~сформированы 


правила перемещения из~$\bm{x}^\prime$ в~целевую точку $\bm{x}\hm = 


(x_1\cdots x_n)$ ($x_j\hm\geq 0$, $j \hm= 1,\ldots , n$): 


  $$


\left\{ F_i(x) = F_i(x^\prime) + h_i[p_i]\right\},


  $$


где $F_i(x) = a_{i1}x_1 +\cdots+ a_{in}x_n$, $i \hm= 1,\ldots , l$; $0 \hm< p_i\hm\leq \infty$ 


для $h_i\not= 0$, $p_i \hm= \infty$ для $h_i = 0$. 


  


  Формально система правил может быть несовместной. Поэтому правила 


с~$h_i \not=0$ рассматриваются как \textit{ориентирующие}, а~$h_i$~--- как 


\textit{задающий шаг} (нередко отличающийся от реального шага, который 


может быть получен для данного набора правил). Если реальные и~задающие 


шаги имеют один и~тот же знак для всех правил, новая точка в~любом случае 


повышает эффективность решения. 


  


  Точка $\bm{x}$ отыскивается следующим образом. Сначала вычисляется 


проекция точки на гиперплоскость 


$$


a_{i1}x_1 +\cdots + a_{in}x_n = F_i(x^\prime) 


+ h_i\  \forall \,h_i\not=0\,.


$$





 Направляющий вектор нормали к~этой гиперплоскости 


равен ($a_{i1}\cdots a_{in}$). Чтобы найти проекцию, необходимо дать переменным 


приращения $ha_{i1}\cdots ha_{in}$ (где $h$~--- неизвестное число). Перемещение 


по нормали дает приращение функции $h(a_{i1}^2 +\cdots + a_{in}^2)$, которое 


должно быть равно~$h$, так что $h \hm= h_i/(a_{i1}^2 +\cdots + a_{in}^2)$ 


(естественно, $a_{i1}^2 +\cdots + a_{in}^2\not= 0$). Таким образом получается 


проекция 


  $$


\left(x_1^\prime + \fr{a_{i1}h_1}{a_{i1}^2 +\cdots+ a_{in}^2}\right)\cdots \left(


x_n^\prime + \fr{a_{in}h_n }{a_{i1}^2 +\cdots + a_{in}^2}\right).


  $$


Когда проекции найдены $\forall \,h_i\not=0$, ищутся требуемые приращения по 


всем переменным: $\Delta x_{ji}\hm = a_{ij}h_i/(a_{i1}^2 +\cdots + a_{in}^2)$ ($j \hm= 


1,\ldots , n$; индексы функций пробегают значения от~1 до~$s$, $1 \hm\leq s 


\hm\leq l$). Далее вычисляется конец среднедействующей  


век\-то\-ров-нор\-малей 


  $$


x_j = x_j^\prime + \fr{p_1\Delta x_{j1} + \cdots + p_s\Delta x_{js}}{p_1 +\cdots + p_s}\quad   


(j = 1,\ldots , n).


  $$


Такая точка ближе к~тем гиперплоскостям, которые соответствуют правилам 


с~более высокими приоритетами. Если приоритеты не были установлены, они 


полагаются равными~1 и~получаются следующие формулы: 


  $$


  x_j=  x_j^\prime + \fr{\Delta x_{j1} + \cdots + \Delta x_{js}}{s}\quad  (j = 1,\ldots , n).


  $$


  


  Если $\exists\,k: h_k = 0$ ($1 \hm\leq k \hm\leq l$), конец среднедействующей 


проецируется на гиперплоскость $a_{k1}x_1 + \cdots + a_{kn}x_n \hm= 


F_k(x^\prime)$; в~противном случае на роль искомой точки~$\bm{x}$ 


претендует конец среднедействующей. Затем контролируется 


неотрицательность переменных ($x_j \hm\geq 0$, $j \hm= 1,\ldots , n$): 


отрицательные обнуляются. И,~наконец, проверяется совпадение знаков 


реальных и~за\-да\-ющих шагов. Если они совпадают для всех правил, 


вычисленная точ\-ка считается целевой точ\-кой~$\bm{x}$. Если не совпадают, 


эксперту предлагается скорректировать сис\-те\-му правил.


  


  Сравнение результатов, полученных с~помощью программно реализованного 


метода целевого перемещения решения, с~результатами, полученными с~по\-мощью программы LINGO~\cite{29-il}, реализующей  


симп\-лекс-ме\-тод~\cite{17-il}, приведено в~\cite{24-il}. 





\section{Заключение}


  


  Подход к~решению линейных задач ситуационного управ\-ле\-ния методом 


целевого перемещения решения создан на основе концепций ситуационного 


управ\-ле\-ния~\cite{13-il, 14-il}, поддержки принятия решений в~нештатных 


ситуациях~\cite{15-il} и~ситуационной информатизации~\cite{16-il}. 


Обновленная версия метода целевого перемещения решения содержит ряд 


уточнений предыдущего варианта~\cite{22-il, 24-il}. 


  


  Технология решения последовательности задач в~режиме вычислительного 


эксперимента ориентирована на реализацию в~виде 


 ин\-тер\-нет-сер\-ви\-са~\cite{12-il}. Эта технология позволяет эксперту (или 


обучаемому роботу) отыскивать решения, соответствующие его 


представлениям о~реализуемости и~эф\-фек\-тив\-ности. Представления, зависящие 


от опыта и~информированности эксперта, могут изменяться с~изменением 


ситуации. 


  


  Программно реализованная технология прошла проверку на задачах 


производственного планирования (по назначению аналогичных 


рас\-смат\-ри\-ва\-емым в~\cite{25-il, 26-il, 27-il, 28-il}), ресурсного обоснования 


решений при планировании операций в~условиях чрезвычайных ситуаций 


и~др.~\cite{24-il}.





{\small\frenchspacing


{%\baselineskip=10.8pt


\addcontentsline{toc}{section}{Литература}


\begin{thebibliography}{99}


\bibitem{1-il}


\Au{Tapscott D.} The digital economy: Promise and peril in the age of networked intelligence.~--- 


New York, NY, USA: McGraw-Hill, 1996. 342~p.


\bibitem{2-il}


\Au{Christensen C.\,M.} The innovator`s dilemma: When new technologies cause great firms to 


fail.~--- Boston, MA, USA: Harvard Business School Press, 1997. {\sf 


http://www. hbs.edu/faculty/Pages/item.aspx?num=46}.


\bibitem{3-il}


 The new digital economy: How it will transform business.~--- Oxford, U.K.: Oxford Economics, 2015. 34~р.


 {\sf http://www.pwc.com/mt/en/publications/assets/the-new-digital-economy.pdf}.


\bibitem{4-il}


Цифровая экономика Российской Федерации: Программа, утвержденная распоряжением 


Правительства Российской Федерации от 28~июля 2017~г.\ №\,1632-р. 87~с. {\sf  


http://d-russia.ru/wp-content/uploads/2017/07/programma-tsifrov-econ.pdf}.





\bibitem{6-il} %5


\Au{Kim R.\,Y.} Efficient wireless communications schemes for machine to machine 


communications~/\!\!/ Comm. Com. Inf. Sc., 2011. Vol.~181. No.\,3. P.~313--323.





\bibitem{5-il} %6


\Au{Lien S.\,Y., Liau T.\,H., Kao~C.\,Y., Chen~K.\,C.} Cooperative access class barring for 


machine-to-machine communications~/\!\!/ IEEE T. Wirel. Commun., 2012. Vol.~11. No.\,1. 


P.~27--32.


\bibitem{7-il} %7


\Au{Perera C., Liu~C.\,H., Jayawardena~S.} The emerging Internet of Things marketplace from 


an industrial perspective: A~survey~/\!\!/ IEEE T. Emerging Topics Computing, 2015. Vol.~3. 


No.\,4. P.~585--598. 


\bibitem{8-il}


<<Интернет вещей>> (IoT) в~России. Технология будущего, доступная уже сейчас~/\!\!/ 


Новости ИТ-канала, 4~августа 2017~г. 64~с.  {\sf https://ru.investinrussia.com/data/file/ IoT-inRussia-research\_rus.pdf}.





\bibitem{9-il}


\Au{Jede A., Teuteberg F.} Understanding socio-technical impacts arising from software  


as-a-service usage in companies~/\!\!/ Bus.  Inf. Syst. Eng., 2016. Vol.~58. 


No.\,3. P.~161--176.





\bibitem{12-il} %10


\Au{Ильин А.\,В.} Интернет-сервисы планирования ресурсов, 2016. {\sf  


https://www.res-plan.com}.





\bibitem{10-il} %11


Five benefits of Software as a~Service. Trumba Corporation, 2017. {\sf 


http://www.trumba.\linebreak com/connect/knowledgecenter/pdf/Saas\_paper\_WP-001.pdf}.


\bibitem{11-il} %12


Internet Services~/\!\!/ Tutorialspoint, 2018. {\sf 


http://www.tutorialspoint.com/internet\_\linebreak technologies/internet\_services.htm}.





\bibitem{13-il}


\Au{Клыков Ю.\,И.} Ситуационное управление большими сис\-те\-ма\-ми.~--- М.: Энергия, 1974. 


241~с.


\bibitem{14-il}


\Au{Поспелов Д.\,А.} Ситуационное управление: тео\-рия и~практика.~--- М.: Наука, 1986. 


288~с.








\bibitem{16-il} %15


\Au{Ильин В.\,Д.} Основания ситуационной информатизации.~--- М.: Наука, Физматлит, 


1996. 180~с.





\bibitem{15-il} %16


\Au{Геловани В.\,А., Башлыков~А.\,А., Бритков~В.\,Б., Вязилов~Е.\,Д.} Интеллектуальные 


системы поддержки принятия решений в~нештатных ситуациях с~использованием 


информации о~со\-сто\-янии природной среды.~--- М.: Эдиториал УРСС, 2001. 304~с.





\bibitem{17-il}


\Au{Dantzig G.\,B.} Linear programming and extensions.~--- Princeton, NJ, USA: Princeton University Press 


and the RAND Corp., 1963. 656~p.


\bibitem{18-il}


\Au{Karmarkar N}. A~new polynomial-time algorithm for linear programming~/\!\!/ 


Combinatorica, 1984. Vol.~4. No.\,4. P.~373--395. 


\bibitem{19-il}


\Au{Tikhonov A.\,N., Arsenin V.\,Y.} Solutions of ill-posed problems.~--- New York, NY, USA: Winston, 1977. 258~p.


\bibitem{20-il}


\Au{Ильин А.\,В., Ильин В.\,Д.} Архитектура вы\-чис\-ли\-тель\-но\-го ядра комплекса про\-граммных 


средств ресурсного обоснования решений.~--- М.: ИПИ РАН, 1995. 23~с.


\bibitem{21-il}


\Au{Ильин В.\,Д., Гавриленко Ю.\,В., Ильин~А.\,В., Макаров~Е.\,М.} Математические 


средства ситуационной информатизации.~--- М.: Наука, Физматлит, 1996, 88~с.


\bibitem{22-il}


\Au{Ильин А.\,В.} Математическое обеспечение процессов преобразования ресурсов~/\!\!/ 


Сис\-те\-мы и~средства информатики, 1999. Вып.~9. С.~159--177.


\bibitem{23-il}


\Au{Ильин А.\,В., Ильин В.\,Д.} Интерактивный преобразователь ресурсов с~изменяемыми 


правилами поведения~/\!\!/ Информационные технологии и~вы\-чис\-ли\-тель\-ные сис\-те\-мы, 2004. 


№\,2. С.~67--77.


\bibitem{24-il}


\Au{Ильин А.\,В.} Экспертное планирование ресурсов.~--- М.: ИПИ РАН, 2013. 58~с.





\bibitem{29-il} %25


\Au{Schrage L.} Optimization modeling with LINGO.~--- Lindo Systems Inc., 1998. 550~p. {\sf 


https://www.lindo.com/downloads/Lingo\_Textbook\_5thEdition.pdf}.


\bibitem{25-il} %26


\Au{Al-Mashari M.} Enterprise resource planning (ERP) systems: A~research agenda~/\!\!/ 


Ind. Manage.  Data Syst., 2002. Vol.~103. No.\,1. P.~22--27. 


\bibitem{26-il} %27


\Au{Umble E.\,J., Haft R.\,R., Umble~M.\,M.} Enterprise resource planning: Implementation 


procedures and critical success factors~/\!\!/ Eur. J.~ Oper. Res., 2003. Vol.~146. No.\,2. 


P.~241--257.


\bibitem{27-il} %28


\Au{Marnewick C., Labuschagne L.} A~conceptual model for enterprise resource planning 


(ERP)~/\!\!/ Information Management Computer Security, 2005. Vol.~13. No.\,2. P.~144--155. 


\bibitem{28-il} %29


\Au{Amoako-Gyampah K.} Perceived usefulness, user involvement and behavioral intention: An 


empirical study of ERP implementation~/\!\!/ Comput. Hum. Behav., 2007. Vol.~23. No.\,3. 


P.~1232--1248. 





 \end{thebibliography}


} }











\vspace*{-6pt}





\hfill{\small\textit{Поступила в~редакцию 14.08.22}}





%\vspace*{8pt}





%\hrule





%\vspace*{2pt}





\newpage





\vspace*{-28pt}





%\hrule





%\vspace*{8pt}





\def\tit{METHOD OF TARGET DISPLACEMENT OF~SOLUTION IN~SITUATIONAL 


MANAGEMENT TECHNOLOGIES}











\def\titkol{Method of target displacement of~solution in~situational 


management technologies}





\def\aut{A.\,V.~Ilyin$^1$ and V.\,D.~Ilyin$^2$}





\def\autkol{A.\,V.~Ilyin and V.\,D.~Ilyin}





\titel{\tit}{\aut}{\autkol}{\titkol}





\vspace*{-8pt}





\noindent


$^1$State Research Institute of Aviation Systems, 7~Viktorenko Str., Moscow 125319,\linebreak


$\hphantom{^1}$Russian  Federation





\noindent


$^2$Federal Research Center ``Computer Science and Control'' of the Russian Academy\linebreak


$\hphantom{^1}$of Sciences, 44-2~Vavilov Str., Moscow 119333, Russian Federation








\def\leftfootline{\small{\textbf{\thepage}


\hfill Sistemy i Sredstva Informatiki~---


Systems and Means of Informatics 2022 vol~32 no\,4}


}%


 \def\rightfootline{\small{Sistemy i Sredstva Informatiki~---


 Systems and Means of Informatics 2022 vol~32 no\,4


\hfill \textbf{\thepage}}}





\vspace*{3pt}


  








\Abste{The updated technology for interactive solving of linear problems of situational 


management by the method of target displacement of solution is considered. The formulations of 


tasks and the requirements for the quality of solutions are based on the analysis of portraits of the 


target and achieved situations. The method of target displacement of solution (as opposed to the 


standard linear programming methods) allows the expert to get the desired result even if the 


constraints system is incompatible. The interactivity of the method allows the expert to change the 


problem formulation at each step of the solution search based on his\!/\!her opinion on the feasibility and 


effectiveness. The technology is focused on implementation in the form of online service.}





\KWE{situational management; linear problems of situational management; resources allocation; 


portrait of the situation; method of target displacement of solution; online service}





\DOI{10.14357\!/\!08696527220412}  





%\vspace*{-10pt}








%\Ack





%\vspace*{-4pt}





%\noindent


%The results were obtained in the course of the research work ``Modeling of social, 


% economic, and ecological processes'' (No.\,0063-2016-0005) carried out in accordance


%  with the state assignment of FANO of Russia for the Federal Research Center 


%  ``Computer Science and Control'' of the Russian Academy of Sciences.














\vspace*{-4pt}





\renewcommand{\bibname}{\protect\rmfamily References}


%\renewcommand{\bibname}{\large\protect\rm References}





{\small\frenchspacing


{ %\baselineskip=12pt


\addcontentsline{toc}{section}{References}


\begin{thebibliography}{99}





\vspace*{-2pt}





\bibitem{1-il-1}


\Aue{Tapscott, D.} 1996. \textit{The digital economy: Promise and peril in the age of networked 


intelligence}. New York, NY: McGraw-Hill. 342~p.


\bibitem{2-il-1}


\Aue{Christensen, C.\,M.} 1997. \textit{The innovator's dilemma: When new technologies cause 


great firms to fail}. Boston, MA: Harvard Business School Press. Available at: {\sf 


http:// www.hbs. edu/faculty/Pages/item.aspx?num=46} (accessed February~4, 2018).


\bibitem{3-il-1}


\textit{The new digital economy: How it will transform business}. 2015. 


Oxford, U.K.: Oxford Economics. 34~p. Available at: 


{\sf http://www.pwc.com/mt/en/publications/assets/the-new-digital-economy.pdf} (accessed 


September~12, 2022).


\bibitem{4-il-1}


 Tsifrovaya ekonomika Rossiyskoy Fe\-de\-ra\-tsii: 


Programma, utverzhdennaya ras\-po\-rya\-zhe\-ni\-em Pravitel'stva Rossiyskoy Federatsii ot 28~iyulya 


2017~g. No.\,\mbox{1632-r} [The program ``Digital Economy of the Russian Federation'' approved by 


Government Order No.\,\mbox{1632-r} dated July~28, 2017]. 87~p. Available at: {\sf  


http://d-russia.ru/wp-content/\linebreak uploads/2017/07/programma-tsifrov-econ.pdf} (accessed 


September~12, 2022).





\bibitem{6-il-1} %5


\Aue{Kim, R.\,Y.} 2011. Efficient wireless communications schemes for machine to machine 


communications. \textit{Comm. Com. Inf. Sc.} 181(3):313--323.





\bibitem{5-il-1} %6


\Aue{Lien, S.\,Y., T.\,H. Liau, C.\,Y. Kao, and K.\,C.~Chen.} 2012. Cooperative access class 


barring for machine-to-machine communications. \textit{IEEE T.~Wirel. Commun.} 11(1):27--32.





\bibitem{7-il-1}


\Aue{Pereyra, C., C.\,H.~Liu, and S.~Jayawardena.} 2015. The emerging Internet of Things 


marketplace from an industrial perspective: A~survey. \textit{IEEE T. Emerging Topics 


Computing} 3(4):585--598.


\bibitem{8-il-1}


PWC. 2017. ``Internet veshchey'' (IoT) v~Rossii. Tekhnologiya budushchego, dostupnaya uzhe 


seychas [``Internet of Things'' (IoT) in Russia. Technology of the future, available now]. 64~p. Available 


at: {\sf https://ru.investinrussia.com/data/file/IoT-inRussia-research\_rus.} (accessed 


October 27, 2022).


\bibitem{9-il-1}


\Aue{Jede, A., and F. Teuteberg.} 2016. Understanding socio-technical impacts arising from 


software as-a-service usage in companies. \textit{Bus. Inf. Syst. Eng.} 58(3):161--176.





\bibitem{12-il-1} %10


\Aue{Ilyin, A.\,V.} 2016. Internet-servisy planirovaniya resursov [Online resource planning 


services]. Available at: {\sf https://www.res-plan.com} (accessed September~12, 2022).





\bibitem{10-il-1} %11


Trumba Corp. 2017. Five benefits of software as a service.  {\sf 


https://www.trumba.\linebreak com/connect/knowledgecenter/pdf/Saas\_paper\_WP-001.pdf} (accessed 


September~12, 2022).


\bibitem{11-il-1} %12


Internet services. 2018. Tutorialspoint. {\sf http://www.tutorialspoint.com/internet\_\linebreak technologies/internet\_services.htm} 


(accessed September 12, 2022).





\bibitem{13-il-1}


\Aue{Klykov, Yu.\,I.} 1974. \textit{Situatsionnoe upravlenie bol'shimi sistemami} [Situational 


management of large systems]. Moscow: Energiya. 241~p.


\bibitem{14-il-1}


\Aue{Pospelov, D.\,A.} 1986. \textit{Situatsionnoe upravlenie: teoriya i~praktika} [Situational 


management: Theory and practice]. Moscow: Nauka. 288~p.





\bibitem{16-il-1} %15


\Aue{Ilyin, V.\,D.} 1996. \textit{Osnovaniya situatsionnoy informatizatsii} [Fundamentals of 


situational informatization]. Moscow: Nauka, Fizmatlit. 180~p.





\bibitem{15-il-1} %16


\Aue{Gelovani, V.\,A., A.\,A.~Bashlykov, V.\,B.~Britkov, and E.\,D.~Vyazilov.} 2001. 


\textit{Intellektual'nye sistemy podderzhki prinyatiya resheniy v~neshtatnykh situatsiyakh 


s~ispol'zovaniem informatsii o~sostoyanii prirodnoy sredy} [Intelligent decision support systems 


in emergency situations using information about the state of the natural environment]. Moscow: 


URSS. 304~p.





\bibitem{17-il-1}


\Aue{Dantzig, G.\,B.} 1963. \textit{Linear programming and extensions}. Princeton, NJ: Princeton 


University Press and the RAND Corp. 656~p.


\bibitem{18-il-1}


\Aue{Karmarkar, N.} 1984. A~new polynomial-time algorithm for linear programming. 


\textit{Combinatorica} 4(4):373--395.


\bibitem{19-il-1}


\Aue{Tikhonov, A.\,N., and V.\,Y.~Arsenin.} 1977. \textit{Solutions of ill-posed problems}. New 


York, NY: Winston. 258~p.


\bibitem{20-il-1}


\Aue{Ilyin, A.\,V., and V.\,D.~Ilyin.} 1995. \textit{Arkhitektura vychislitel'nogo yadra kompleksa 


programmnykh sredstv resursnogo obosnovaniya resheniy} [Architecture of computational kernel 


of software designed to search for resource management decisions]. Moscow: Institute of 


Informatics Problems of the Russian Academy of Sciences. 23~p.


\bibitem{21-il-1}


\Aue{Ilyin, V.\,D., Yu.\,V.~Gavrilenko, A.\,V.~Il'in, and E.\,M.~Makarov.} 1996. 


\textit{Ma\-te\-ma\-ti\-che\-skie sredstva situatsionnoy informatizatsii} [Mathematical means of situational 


informatization]. Moscow: Nauka, Fizmatlit. 88~p.


\bibitem{22-il-1}


\Aue{Ilyin, A.\,V.} 1999. Matematicheskoe obespechenie protsessov preobrazovaniya re\-sur\-sov 


[Mathematical means of resource transformation processes]. \textit{Sistemy i~Sredstva 


Informatiki~--- Systems and Means of Informatics} 9:159--177.


\bibitem{23-il-1}


\Aue{Ilyin, A.\,V., and V.\,D.~Ilyin.} 2004. Interaktivnyy preobrazovatel' resursov 


s~izmenyaemymi pravilami povedeniya [Interactive resource converter with customisable rules of 


behavior]. \textit{Informatsionnye tekhnologii i~vychislitel'nye sistemy} [J.~Information Technologies 


and Computing Systems] 2:67--77.


\bibitem{24-il-1}


\Aue{Ilyin, A.\,V.} 2013. \textit{Ekspertnoe planirovanie resursov} [Expert resource planning]. 


Moscow: Institute of Informatics Problems of the Russian Academy of Sciences. 58~p.





\bibitem{29-il-1} %25


\Aue{Schrage, L.} 1998. \textit{Optimization modeling with LINGO}. Lindo Systems Inc. 550~p. 


{\sf 


https://www.lindo.com/downloads/Lingo\_Textbook\_5thEdition.pdf}.


\bibitem{25-il-1} %26


\Aue{Al-Mashari, M.} 2002. Enterprise resource planning (ERP) systems: A research agenda. 


\textit{Ind. Manage. Data Syst.} 103(1):22--27. 


\bibitem{26-il-1} %27


\Aue{Umble, E.\,J., R.\,R. Haft, and M.\,M.~Umble.} 2003. Enterprise resource planning: 


Implementation procedures and critical success factors. \textit{Eur. J.~Oper. Res.}  


146(2):241--257.


\bibitem{27-il-1} %28


\Aue{Marnewick, C., and L. Labuschagne.} 2005. A~conceptual model for enterprise resource 


planning (ERP). \textit{Information Management Computer Security} 13(2):144--155.


\bibitem{28-il-1} %29


\Aue{Amoako-Gyampah, K.} 2007. Perceived usefulness, user involvement and behavioral 


intention: An empirical study of ERP implementation. \textit{Comput. Hum. Behav.}  


23(3):1232--1248. 





\end{thebibliography}


} }





\vspace*{-6pt}





\hfill{\small\textit{Received August 14, 2022}} 





%\vspace*{-8pt} 





\Contr





\noindent


\textbf{Ilyin Alexander V.} (b.\ 1975)~--- Candidate of Science (PhD) in technology, leading 


scientist, State Research Institute of Aviation Systems, 7~Viktorenko Str., Moscow 125319, 


Russian Federation; \mbox{ilyin@res-plan.com}





\vspace*{3pt}





\noindent


\textbf{Ilyin Vladimir D.} (b.\ 1937)~--- Doctor of Science in technology, professor, leading 


scientist, A.\,A.~Dorodnicyn Computing Center, Federal Research Center ``Computer Science and 


Control'' of the Russian Academy of Sciences, 40~Vavilov Str., Moscow119333, Russian 


Federation; \mbox{vdilyin@yandex.ru}


  


\label{end\stat}





\renewcommand{\bibname}{\protect\rm Литература} 


